%-------------------------------------------------------------------------------
%	SECTION TITLE
%-------------------------------------------------------------------------------
\cvsection{Summary}


%-------------------------------------------------------------------------------
%	CONTENT
%-------------------------------------------------------------------------------
\begin{cvparagraph}

%---------------------------------------------------------
I am a computational scientist and AI researcher developing machine learning and statistical methods for biomedical discovery. I am currently a Postdoctoral Research Fellow at the Cancer Data Science Laboratory (CDSL), NCI/CCR, NIH, where my research focuses on computational immuno-oncology and the inference of cell--cell communication. In particular, I develop methods to infer the activity of secreted proteins and characterize how signaling is coordinated across cells and tissues using bulk, single-cell, and spatial transcriptomics. This work integrates biological knowledge, statistical modeling, and machine learning to uncover signaling mechanisms, identify clinically relevant biomarkers, and improve our understanding of therapeutic response.

My research has led to computational frameworks and software for biological activity inference, including SecAct, a framework for inferring secreted-protein signaling activities across more than 1,000 human secreted proteins, and its Python implementation, as well as GPU-accelerated tools for large-scale spatial analysis. I have also developed machine learning approaches for clinical biomarker discovery, including IC2Bert for immune checkpoint blockade response prediction, and contributed to broader studies of artificial intelligence in precision oncology. My work has been published in \textit{Nature Methods}, \textit{Nature Cancer}, \textit{Scientific Reports}, and other peer-reviewed journals, and has resulted in open-source software and three issued patents.

More broadly, my methodological interests include biologically informed machine learning, representation learning, foundation and self-supervised learning, sparse modeling, and statistical methods for evaluating biological data and model quality. Before focusing on computational immuno-oncology, I worked on gene-expression-based disease and drug-response prediction, biosignal processing and diagnostic modeling, and microfluidic systems for synthetic biology. This interdisciplinary training, spanning mechanical engineering, bioengineering, computational biology, and AI, allows me to approach biomedical problems from both methodological and biological perspectives and to develop computational approaches that translate across data modalities and research domains.

\end{cvparagraph}
